\documentclass[12pt,a4paper]{article}

\usepackage[utf8]{inputenc}
\usepackage[T5]{fontenc}
\usepackage[utf8]{vntex}

\usepackage{lmodern}
\usepackage{geometry}
\geometry{margin=2.5cm}
\usepackage{graphicx}
\graphicspath{{figures/}{./}}
\usepackage{amsmath}
\usepackage{booktabs}
\usepackage{float}
\usepackage{array}
\setlength{\parindent}{0pt}
\setlength{\parskip}{6pt}
\usepackage{caption}
\usepackage{subcaption}
\usepackage{xcolor}
\usepackage[hypertexnames=false]{hyperref}
\usepackage{longtable}

% English names. vntex installs Vietnamese caption names from a
% begin-document hook, so ours must be deferred to the same hook to win.
\AtBeginDocument{%
  \renewcommand{\contentsname}{Contents}%
  \renewcommand{\listfigurename}{List of Figures}%
  \renewcommand{\listtablename}{List of Tables}%
  \renewcommand{\figurename}{Figure}%
  \renewcommand{\tablename}{Table}%
  \renewcommand{\refname}{References}%
}

% Keep floats inside their section
\usepackage[section]{placeins}

\hypersetup{
    colorlinks=true,
    linkcolor=black,
    urlcolor=blue,
}

\setlength{\floatsep}{10pt plus 2pt minus 2pt}
\setlength{\textfloatsep}{12pt plus 2pt minus 4pt}
\setlength{\intextsep}{10pt plus 2pt minus 2pt}
\renewcommand{\textfraction}{0.05}
\renewcommand{\floatpagefraction}{0.85}
\setlength{\emergencystretch}{3em}
\hbadness=2000

% Figure helper: insert the real image when the file is present, otherwise a
% clean placeholder box, so the report compiles before the diagrams exist.
% Drop the named .png into this folder later and the figure upgrades itself.
\newcommand{\reportfig}[3]{%
  \IfFileExists{figures/#2}{%
    \includegraphics[width=#1\textwidth]{#2}%
  }{%
    \fbox{\begin{minipage}[c][5.5cm][c]{#1\textwidth}%
      \centering\itshape
      [\,Figure to be inserted\,]\par\vspace{6pt}
      #3\par\vspace{10pt}%
      \normalfont\ttfamily\footnotesize #2%
    \end{minipage}}%
  }%
}

% Logo helper: same idea for the title-page logos, so a missing logo file does
% not break compilation.
\newcommand{\titlelogo}[1]{%
  \IfFileExists{#1}{\includegraphics[width=0.9\textwidth]{#1}}{%
    \fbox{\parbox[c][2.6cm][c]{0.8\textwidth}{\centering\ttfamily\footnotesize #1}}%
  }%
}

\begin{document}
\raggedbottom
\pagenumbering{gobble}

% =========================================================================
% TITLE PAGE
% =========================================================================
\begin{titlepage}
    \centering
    \textbf{\large University of Engineering and Technology, VNU, Hanoi}\\[0.2cm]
    \textbf{\large Institute for Artificial Intelligence}\\[0.5cm]
    \rule{10cm}{0.4pt}\\

    \vspace{2cm}
    \begin{minipage}{0.45\textwidth}
        \begin{center}
            \titlelogo{uet_logo.jpg}
        \end{center}
    \end{minipage}
    \hfill
    \begin{minipage}{0.45\textwidth}
        \begin{center}
            \titlelogo{ai_logo.png}
        \end{center}
    \end{minipage}

    \vspace{1.5cm}
    {\large \textbf{COURSE REPORT}}\\[0.5cm]
    {\large \textbf{INDUSTRIAL INTERNSHIP}}\\[0.5cm]

    \vspace{5pt}
    {\large{
    Topic: \textbf{HumJob: A Local-First, Grid-Aware Pipeline for
    Transcribing a Hummed Melody into Notes, Key and Chords}
    }}\\[0.25cm]

    \vspace{12pt}
    {\large
    \begin{tabular}{>{\bfseries}l l}
        Supervisor: & PhD. Trần Hồng Việt \\
    \end{tabular}
    }\\[0.5cm]

    \vspace{12pt}
    {\large
    \begin{tabular}{>{\bfseries}l l}
        Author: & Vũ Nguyên Đan - 23020351 \\
    \end{tabular}
    }\\[0.5cm]

    \vfill
    \centering
    \textbf{Hà Nội, 2026}
\end{titlepage}

% =========================================================================
% EVALUATION PAGE
% =========================================================================
\newpage
\begin{center}
    \textbf{Supervisor's Evaluation:}
\end{center}

\vspace{0.2cm}

\noindent\makebox[\linewidth]{\dotfill}\\[0.2cm]
\noindent\makebox[\linewidth]{\dotfill}\\[0.2cm]
\noindent\makebox[\linewidth]{\dotfill}\\[0.2cm]
\noindent\makebox[\linewidth]{\dotfill}\\[0.2cm]
\noindent\makebox[\linewidth]{\dotfill}\\[0.2cm]
\noindent\makebox[\linewidth]{\dotfill}\\[0.2cm]
\noindent\makebox[\linewidth]{\dotfill}\\[0.2cm]
\noindent\makebox[\linewidth]{\dotfill}\\[0.2cm]
\noindent\makebox[\linewidth]{\dotfill}\\[0.2cm]
\noindent\makebox[\linewidth]{\dotfill}\\[0.2cm]
\noindent\makebox[\linewidth]{\dotfill}\\[0.2cm]
\noindent\makebox[\linewidth]{\dotfill}\\[0.2cm]
\noindent\makebox[\linewidth]{\dotfill}\\[0.2cm]
\noindent\makebox[\linewidth]{\dotfill}\\[0.2cm]
\noindent\makebox[\linewidth]{\dotfill}\\[0.2cm]
\noindent\makebox[\linewidth]{\dotfill}\\[0.2cm]
\noindent\makebox[\linewidth]{\dotfill}\\[0.2cm]
\noindent\makebox[\linewidth]{\dotfill}

\noindent\hspace{3cm} Score (in figures) \dotfill{} Score (in words) \dotfill

\vspace{1.5cm}

\begin{center}
    Hà Nội, \ldots{} \ldots{} 2026\\[0.2cm]
    \textbf{Supervisor}\\[0.1cm]
    (Signature and full name)
\end{center}

% =========================================================================
% ACKNOWLEDGEMENTS
% =========================================================================
\newpage
\begin{center}
    \textbf{\large ACKNOWLEDGEMENTS}
\end{center}

I would like to thank my supervisor for the guidance and feedback received
throughout the Industrial Internship course, which shaped both the direction of this work,
\textbf{HumJob}, and the standard of intellectual honesty I have tried to hold
it to.

This report is the work of a single student. Given the breadth of the system
and the limited time available, it inevitably contains shortcomings. In
particular, the central finding of the evaluation is a negative one: the
synthetic benchmark the project was tuned against no longer discriminates,
while the system is not yet reliable on real hummed input. I have chosen to
report that honestly rather than to present a saturated benchmark as a
success. I welcome any comments and corrections from the instructor so that
the work may be improved.

% =========================================================================
% ABSTRACT
% =========================================================================
\newpage
\begin{center}
    \textbf{\large ABSTRACT}
\end{center}

\textbf{HumJob} is a local-first system that transcribes a hummed melody into
notation. It returns the note sequence, the estimated musical key, and a set
of suggested chords, and exports them as MIDI, MusicXML, and engraved sheet
music. The underlying problem is monophonic melody transcription from a bare
human voice, which is hard because a continuous pitch contour has to be cut
into discrete notes with correct pitches, onsets, and durations.

HumJob makes the problem tractable with two deliberate user-interface
constraints. The user hums the syllable ``da-da-da'', whose consonant gives
every note a crisp onset and separates repeated notes of the same pitch; and
the user hums to a metronome at a known tempo, which turns rhythm from
estimation into snapping onto a known grid. The system is built as a chain of
pure functions over three data types, with every tunable parameter
centralised, and supports five interchangeable pitch backends (pyin, CREPE,
PESTO, FCNF0++, and basic-pitch). A grid-aware segmenter and a
backend-agnostic consolidation stage address the dominant failure mode, and a
spectral octave-correction stage repairs a systematic pitch-tracker error on
legato lines.

The system is evaluated on a synthetic fixture harness that scores note
detection (pitch and onset F1) and, separately, rhythm (quantised onset and
duration accuracy against an intended grid). On the expressive ``realistic''
fixtures the note F1 rose from $0.799$ to $1.000$ across three engineering
iterations, and the harness is now saturated. This report treats that
saturation as a limitation rather than a result: the synthetic benchmark no
longer discriminates, informal testing on real hums still fails, and no
labelled dataset of real hummed melodies exists to measure against. The rhythm
evaluation isolates the concrete mechanism behind the real-world failures.
Quantisation is perfect on-grid but degrades under human timing jitter (mean
rhythm accuracy $0.937$) and collapses under a $5\%$ tempo error ($0.602$), and
the tempo-detection component that is meant to prevent that error is a single
global-tempo estimator that is fragile on real, variable-tempo humming. This
report documents the architecture, the evaluation, and an honest account of
what a saturated synthetic benchmark can and cannot demonstrate, together with
a concrete plan for obtaining real ground truth.

% =========================================================================
% TABLE OF CONTENTS
% =========================================================================
\newpage
\pagenumbering{arabic}
\tableofcontents
\newpage
\listoffigures
\newpage
\listoftables
\newpage

% =========================================================================
% SECTION 1: INTRODUCTION
% =========================================================================
\section{Introduction}

\subsection{Motivation and significance}

A common musical situation has no good tool. A person has a melody in their
head, hums it, and wants it written down: as notes on a staff, as a MIDI file
they can edit, or as a chord sketch they can play. Turning that hum into
notation is the task known as \emph{query-by-humming} transcription, and it is
harder than it looks. The human voice produces a single continuous pitch
contour that slides, wavers, and drifts; notation requires that contour to be
cut into a sequence of discrete notes, each with a definite pitch, a definite
start time, and a definite written duration, and it requires a key and a
metrical grid within which those durations are meaningful.

Two families of existing tools sidestep parts of the problem rather than
solving it. General audio-to-MIDI transcribers are trained on instruments and
polyphonic music; on a bare voice they are prone to octave jumps and they do
not exploit any of the structure a deliberate hum provides. Real-time pitch
monitors and tuners report the instantaneous pitch but never commit to a note
segmentation at all. What is missing is a system that is built specifically for
a cooperative human who is willing to hum in a way that makes transcription
tractable.

HumJob is that system. Its central design decision is to accept two small
constraints on how the user hums, and to build the entire pipeline around them:

\begin{itemize}
    \item The user hums the syllable \textbf{``da-da-da''}. The consonant
    closes the vocal tract briefly at the start of each note, which produces a
    crisp amplitude onset and, crucially, separates two repeated notes of the
    same pitch that a legato hum would smear into one.
    \item The user hums \textbf{to a metronome at a known tempo}. A known
    tempo grid converts rhythm from a hard estimation problem into a snapping
    problem: note onsets and durations are aligned to the nearest grid
    position rather than guessed from raw timing.
\end{itemize}

These constraints are load-bearing. They are what make the problem tractable
for a single-student project, and a recurring theme of this report is that the
remaining failures of the system are failures to exploit them fully, not
failures of the underlying pitch models.

\subsection{Objectives}

The project pursues the following objectives:

\begin{itemize}
    \item To design and implement a complete, local-first melody-transcription
    pipeline that takes a hummed recording and returns notes, key, and
    suggested chords, exported as MIDI, MusicXML, and engraved sheet music.
    \item To structure the pipeline as a chain of pure functions over a small,
    fixed data model, so that every stage is independently testable and every
    tunable parameter is centralised rather than scattered as magic numbers.
    \item To support several interchangeable pitch-detection backends behind a
    single interface, so the note-production model can be chosen or replaced
    without touching the rest of the pipeline.
    \item To deliver the pipeline both as a command-line tool and as a browser
    web application that records audio locally and calls the same code, with an
    interactive editor for correcting the automatic draft.
    \item To build an evaluation harness that measures the system honestly, and
    to characterise clearly what that harness can and cannot demonstrate.
\end{itemize}

\subsection{Research questions}

The project is organised around three research questions:

\begin{itemize}
    \item \textbf{RQ1 - Do the interface constraints help?} Do the two
    deliberate constraints (a consonant onset per note, and humming to a known
    tempo grid) turn melody transcription from an estimation problem into a
    tractable segmentation and snapping problem?
    \item \textbf{RQ2 - Is this a segmentation problem or a model problem?}
    When notes come out wrong, does the fault lie chiefly in the segmentation
    and quantisation of an otherwise-correct pitch contour, rather than in the
    pitch model itself?
    \item \textbf{RQ3 - How well can the system be evaluated?} How well does
    the system agree with ground truth, and what are the methodological limits
    of that measurement given that the project has only synthetic ground truth
    and no labelled corpus of real hummed melodies?
\end{itemize}

% =========================================================================
% SECTION 2: BACKGROUND
% =========================================================================
\section{Background}

\subsection{Melody transcription and the monophonic subproblem}

Automatic music transcription in general is the conversion of an audio signal
into a symbolic score. The general problem is polyphonic and very hard. HumJob
addresses the restricted \emph{monophonic} case: one voice, one note at a time.
This restriction removes the need to separate simultaneous pitches, but it does
not make the problem trivial, because the difficulty moves from
source-separation to \emph{segmentation}: deciding where one note ends and the
next begins, and how long each note is in metrical terms.

A monophonic transcription pipeline has three logical stages: estimate a pitch
contour over time, cut that contour into discrete note events, and place those
events on a key and a rhythmic grid. HumJob follows exactly this shape, and the
bulk of its engineering effort is in the second and third stages.

\subsection{Fundamental-frequency estimation}

The first stage estimates the fundamental frequency ($f_0$) of the voice, frame
by frame. Several families of estimator exist and HumJob supports five of them
behind one interface:

\begin{itemize}
    \item \textbf{pYIN}~\cite{pyin} is a classical signal-processing estimator
    that augments the YIN algorithm with probabilistic thresholds and a Viterbi
    decoding step. It is fast, dependency-light, and the project default.
    \item \textbf{CREPE}~\cite{crepe} is a convolutional neural network trained
    to estimate pitch directly from the waveform, and is steady on singing
    voice.
    \item \textbf{PESTO}~\cite{pesto} is a self-supervised,
    transposition-equivariant pitch estimator; it is the web application's
    default because it is precise on voice while being lighter and faster than
    CREPE.
    \item \textbf{FCNF0++}~\cite{penn} is a fully-convolutional pitch and
    periodicity estimator, a precision peer of PESTO.
    \item \textbf{basic-pitch}~\cite{basicpitch} is an instrument-trained
    neural note-transcription model that emits note events directly, bypassing
    the frame-level contour.
\end{itemize}

The important background fact for this report concerns a specific failure of
Viterbi-decoded trackers. On a continuously voiced legato line, with no silence
to reset the decoder, the ``stay put'' transition prior of a Viterbi decoder
can lock a whole note onto its subharmonic, reporting it a full octave too low.
This is not an audio artifact; it is a decoding artifact, and Section~%
\ref{sec:octave} describes how HumJob detects and repairs it.

\subsection{From a pitch contour to notes}

Turning a contour into notes requires three sub-decisions. \emph{Voicing}
decides which frames carry pitch at all, gating out breaths and silence.
\emph{Segmentation} decides where the boundaries between notes fall, which for
a ``da-da-da'' hum means finding the brief energy dips created by the
consonant. \emph{Consolidation} does the opposite where needed: a single held
note that vibrato has fragmented into several pieces must be fused back into
one. Segmentation and consolidation pull in opposite directions, and balancing
them is the core difficulty of the pipeline.

\subsection{Key finding and rhythm quantisation}

Once notes exist, two musical decisions remain. \emph{Key finding} identifies
the scale the melody sits in; HumJob correlates the sung pitch-class
distribution against the Krumhansl-Kessler probe-tone profiles for all
twenty-four major and minor keys~\cite{krumhansl}. \emph{Quantisation} snaps
note onsets and durations to a metrical grid derived from the known tempo,
which is what allows durations to be written as clean note values rather than
as awkward fractions. When the tempo is wrong, quantisation produces
non-integer durations that notation software renders as strings of tied notes,
which is the visible symptom of most rhythm errors.

\subsection{Why a metronome and ``da-da-da''}

The two constraints introduced in Section~1 are the project's answer to the two
hardest sub-problems above. The consonant in ``da-da-da'' directly supplies the
onset cue that segmentation needs, and in particular it is the only reliable
way to separate two consecutive notes of the same pitch. The metronome supplies
the grid that quantisation snaps to. Neither constraint is burdensome for a
cooperative user, and together they replace two estimation problems with two
detection problems. The project treats them as fixed premises, not as
limitations to be optimised away.

% =========================================================================
% SECTION 3: SYSTEM DESIGN AND METHODOLOGY
% =========================================================================
\section{System Design and Methodology}

\subsection{Architecture overview}

HumJob is organised as a chain of pure functions over three data types.
Figure~\ref{fig:arch} shows the high-level structure: a core Python package
implements the pipeline, a command-line tool calls it directly, and a FastAPI
web application records audio in the browser and calls the same pipeline
server-side. Everything runs locally, and no audio is ever sent to any
third-party service. The single exception to local-only operation is an opt-in
coaching feature (Section~\ref{sec:coaching}): only when the user explicitly
presses a button does it send a short numeric summary of a sung practice take,
never any audio, to an external language-model API.

\begin{figure}[H]
\centering
\reportfig{0.9}{architecture.png}{System architecture: the core pipeline
package, the command-line entry point, and the FastAPI web application with its
static browser front-end. The same pipeline code serves both entry points.}
\caption{High-level system architecture of HumJob.}
\label{fig:arch}
\end{figure}

The three data types are the contract that every stage shares
(Table~\ref{tab:types}). A \texttt{Frame} is one dense analysis frame, roughly
one per $11.6$ ms hop, carrying the frame time, the estimated $f_0$, a voicing
confidence, and a short-time energy. A \texttt{NoteEvent} is one discrete note
after segmentation, carrying its start and end in seconds, its integer pitch,
its raw fractional pitch, a cents offset from the nearest semitone, and, once
quantised, its grid-aligned position and duration in quarter-note units. A
\texttt{Score} is the whole transcription, ready to export.

\begin{table}[H]
\centering
\caption{The three data types the pipeline is built on.}
\label{tab:types}
\begin{tabular}{l p{0.60\textwidth}}
\toprule
\textbf{Type} & \textbf{Role} \\
\midrule
\texttt{Frame}     & One analysis frame ($\approx 11.6$ ms): frame time, $f_0$
(NaN if unvoiced), voicing confidence, short-time energy. What the pitch
tracker sees. \\
\texttt{NoteEvent} & One discrete note after segmentation: start/end seconds,
integer MIDI pitch, raw fractional pitch, cents offset, and (after
quantisation) grid position and duration in quarter-note units. \\
\texttt{Score}     & The whole transcription: the note sequence, the detected
key, tuning offset, and chords, ready to export. \\
\bottomrule
\end{tabular}
\end{table}

A design rule the project holds to is that every tunable parameter lives in a
single central parameter object, not as a magic number inside a stage. Changing
behaviour means changing a named knob, which keeps the stages readable and the
experiments reproducible.

\subsection{The transcription pipeline}

Figure~\ref{fig:pipeline} shows the pipeline stages in order. The chain is:

\begin{center}
audio $\rightarrow$ preprocess $\rightarrow$ [note production] $\rightarrow$
consolidate $\rightarrow$ tuning $\rightarrow$ key $\rightarrow$ quantise
$\rightarrow$ chords $\rightarrow$ export.
\end{center}

\begin{figure}[H]
\centering
\reportfig{0.92}{pipeline-transcriber.png}{The transcription pipeline. Note
production has two paths: the per-frame trackers (pyin / CREPE / PESTO /
FCNF0++) feed octave correction, voicing, and the grid-aware segmenter, while
basic-pitch emits note events directly. Both paths then share consolidation,
tuning, key finding, quantisation, chords, and export.}
\caption{The HumJob transcription pipeline, both note-production paths.}
\label{fig:pipeline}
\end{figure}

\textbf{Note production} has two paths, selected by a backend parameter. The
neural note-event backend (basic-pitch) emits notes directly and has no frame
contour. The per-frame trackers (pyin, CREPE, PESTO, FCNF0++) produce a
\texttt{Frame} contour that then passes through octave correction, a voicing
gate, and the grid-aware segmenter to become notes.

\textbf{Consolidation} runs for every backend. It fuses the fragments that a
held, vibrato-laden note leaves behind, which is the fix for the ``one note
becomes many slivers'' failure. It is grid-aware: it will not fuse two
same-pitch notes across a grid onset, so it never undoes a re-articulation that
the segmenter deliberately split on the beat.

\textbf{The remaining stages} apply a global tuning correction, find the key by
the Krumhansl-Kessler correlation, quantise onsets and durations onto the
tempo grid, suggest chords, and export to MIDI, MusicXML, and engraved sheet
music.

\subsection{The note-detection backends}

The five backends trade precision, speed, and training domain differently.
Table~\ref{tab:backends} summarises them. The default differs by entry point,
which is intentional: the library default is pyin (fast and dependency-light),
the command-line default is basic-pitch, and the web application default is
PESTO (most precise on voice). The four per-frame trackers all feed the same
segmenter; only basic-pitch bypasses it.

\begin{table}[H]
\centering
\caption{The note-detection backends.}
\label{tab:backends}
\begin{tabular}{l p{0.34\textwidth} p{0.34\textwidth}}
\toprule
\textbf{Backend} & \textbf{What it is} & \textbf{Best for / notes} \\
\midrule
\texttt{pesto} & Self-supervised, transposition-equivariant pitch estimator &
Most precise on voice; web-app default; lighter and faster than CREPE \\
\texttt{fcnf0} & FCNF0++, a fully-convolutional $f_0$ estimator & Precision
peer of PESTO; shines on real voice; optional install \\
\texttt{crepe} & Neural pitch CNN, voice-trained & Steady on humming and
singing; uses the accurate full model with Viterbi decoding \\
\texttt{basic\_pitch} & Instrument-trained note-transcription CNN & Instrument
clips; can octave-jump on a bare voice; bypasses the segmenter \\
\texttt{pyin} & Classical signal-processing $f_0$ plus the project segmenter &
Crisp staccato ``da-da-da''; library default \\
\bottomrule
\end{tabular}
\end{table}

\subsection{Grid-awareness: segmentation and consolidation}

The single most important methodological idea in the pipeline is that
segmentation is \emph{grid-aware}. The segmenter is passed the known tempo and
a fine short-window energy envelope. A short or soft consonant closure between
two same-pitch notes appears as a narrow dip in that fine envelope; a width
gate separates such a closure from a wide tremolo swell, and a shallower dip is
still accepted when it lands on a beat, because the grid says a note boundary is
expected there. This is what separates repeats of the same pitch, and it
replaced an earlier coarse valley splitter that smeared over short closures.

Consolidation is the necessary counterweight. Because vibrato and breath can
fragment one held note into several pieces, a backend-agnostic consolidation
stage fuses same-pitch fragments back together. Its grid-awareness is the safety
interlock: it refuses to fuse across a grid onset, so segmentation and
consolidation cannot fight each other into an oscillation.

\subsection{Spectral octave correction}
\label{sec:octave}

A distinct stage, inserted between the tracker and the voicing gate for the
per-frame backends, repairs the Viterbi subharmonic error described in
Section~2. The stage is backend-agnostic and rests on a spectral signature. A
genuine fundamental at frequency $f$ has energy at both its odd harmonics
($f, 3f, 5f$) and its even harmonics ($2f, 4f, 6f$). A subharmonic error, in
which the tracker has reported half the true frequency, has energy only at the
even multiples, because those are the true fundamental and its harmonics; the
odd multiples of the reported frequency have little or no energy. The stage
therefore computes, per voiced frame, the salience at the odd and even
multiples of the reported $f_0$, and doubles the reported frequency only when
the odd-harmonic salience collapses relative to the even-harmonic salience. The
test is deliberately conservative: it leaves genuine fundamentals alone,
including missing-fundamental voices, which retain energy at their odd
harmonics. Section~\ref{sec:results-note} reports the effect of this stage on
the evaluation.

\subsection{The interactive editor (Manual mode)}

The automatic draft is not always correct, and the weak link is segmentation
rather than the pitch contour. The web application therefore provides a Manual
mode: an in-browser, staff-based editor for fixing the automatic draft. It is
almost entirely client-side. It re-implements the score-building and
enharmonic-spelling logic of the export stage in the browser, engraves with a
vendored WebAssembly build of the Verovio music-notation
engine~\cite{verovio}, and offers pure editing operations (change pitch, change
duration, merge, split, delete to rest, insert) with undo and redo. A reference
pitch strip shows the hummed contour against the chosen notes, so the user can
see where the draft diverged from what they sang. The client sends the edited
notes back to the server only at the edges, to re-export MIDI and MusicXML and
to recompute key and chords.

\subsection{Other independent paths}

The web application is tabbed, and four tabs are separate paths that do not go
through the transcription segmenter:

\begin{itemize}
    \item \textbf{Pitch Finder} analyses an uploaded clip for key, tempo, and a
    Camelot code, using a chroma-based method that works on polyphonic audio.
    It reuses the key-scoring and tempo-folding code but not the monophonic
    segmenter, which would be wrong for a full song.
    \item \textbf{Realtime} is a client-side live pitch monitor and guitar
    tuner built on the Web Audio API, with a set of vocal-training drills (a
    target-note lane, a match game, a scale trainer, vibrato analysis, a vocal
    range finder, and a per-session progress log). A server round-trip cannot
    be real-time, so this path runs entirely in the browser and only asks the
    server for the sung key on stop, from a pitch histogram, with no audio
    upload.
    \item \textbf{Transposer} shifts a score to a new key. A file mode
    transposes an uploaded MIDI or MusicXML file server-side with
    music21~\cite{music21}, moving every voice and the key signature; a hum
    mode transposes the last transcription client-side. The file mode also
    offers Camelot-compatible key presets as one-click shifts.
    \item \textbf{Sing-Along} turns an uploaded MIDI or MusicXML score into a
    karaoke exercise: the melody is played back as a guide while the user's
    voice is tracked live and scored note by note. Because the whole melody is
    known in advance, the score itself drives the clock; there is no free-timing
    alignment. The server reduces a possibly polyphonic upload to a single
    monophonic melody line (its skyline, with ties stripped), and the browser
    schedules the count-in and the guide notes at absolute audio-clock times and
    scores the sung pitch against the active note. The in-tune tolerance is
    selectable (four difficulty levels), the guide volume is adjustable, and the
    take can be paused and resumed.
\end{itemize}

\subsection{Optional AI coaching and the local-first exception}
\label{sec:coaching}

The Sing-Along tab includes an optional coaching feature, and it is the single
place in the system where data leaves the machine. After a take, the client
computes a deterministic, human-readable analysis of the performance entirely
in the browser: the signed pitch bias (whether the singer was consistently
sharp or flat), any drift from the start of the take to the end, octave slips,
the accuracy of leaps versus stepwise motion, the weakest register, and the
worst individual notes with their bar numbers. This analysis is always shown
and requires no network access or key.

Only if the user then presses a ``Get coaching'' button is a compact,
privacy-preserving numeric summary of that analysis, together with the musical
context (key, tempo, time signature, chosen difficulty), sent to an external
large-language-model API, which returns spoken-language feedback and practice
suggestions in the user's chosen language. The request contains numbers only:
no audio, no recording, and not even the file name is transmitted. The API key
is read server-side from a local, git-ignored configuration file and is never
exposed to the browser, and the model's reply is rendered as plain text rather
than as markup, so untrusted model output cannot inject content into the page.
The feature is off by default and does nothing until a key is configured, so the
rest of the system remains strictly local; this one call is the deliberate,
clearly-scoped exception to HumJob's otherwise local-only operation.

\subsection{Tempo detection and its known fragility}
\label{sec:tempo-design}

Because a wrong tempo is the usual cause of tied-note artifacts, the
application offers a ``find my tempo'' button. Its purpose is narrow: the user
hums a few even beats, the system estimates a comfortable tempo, and the user
then records the real take to a click at that tempo, so that quantisation lands
on whole beats.

The estimator computes an onset-strength envelope with
librosa~\cite{librosa}, then finds the dominant periodicity of that envelope by
autocorrelation, weighted by a prior centred on a comfortable humming tempo,
and folds the result into a plausible range by octaves. The important property,
and a known weakness examined in Section~\ref{sec:bpm}, is that this returns a
\emph{single global tempo} averaged over the whole clip, tilted toward the
prior. It assumes a roughly constant tempo and reasonably clean onsets, and it
is fragile when either assumption is violated.

\subsection{Technology stack and running the system}
\label{sec:stack}

Table~\ref{tab:stack} lists the main components. The whole system runs locally.
The web application is a FastAPI server that serves a static browser front-end
and reloads on source edits; the command-line tool transcribes a file directly.

\begin{table}[H]
\centering
\caption{Technology stack.}
\label{tab:stack}
\begin{tabular}{l p{0.60\textwidth}}
\toprule
\textbf{Layer} & \textbf{Technology} \\
\midrule
Core pipeline     & Python 3.12, NumPy \\
Pitch backends    & pYIN via librosa~\cite{librosa}; torchcrepe~\cite{crepe};
pesto-pitch~\cite{pesto}; penn / FCNF0++~\cite{penn}; basic-pitch on ONNX
Runtime~\cite{basicpitch} \\
Onset / tempo / chroma & librosa~\cite{librosa} \\
Score model, MIDI, MusicXML & music21~\cite{music21} \\
Engraving         & Verovio~\cite{verovio} (server-side, and a vendored
WebAssembly build in the browser) \\
Web backend       & FastAPI, Uvicorn (with automatic reload) \\
Web front-end     & Static HTML / CSS / JavaScript, Web Audio API \\
AI coaching (optional) & httpx client to an external, OpenAI-compatible
language-model API (DeepSeek); off by default, opt-in, numeric summary only \\
Evaluation        & mir\_eval~\cite{mireval} for note metrics \\
\bottomrule
\end{tabular}
\end{table}

% =========================================================================
% SECTION 4: EVALUATION AND RESULTS
% =========================================================================
\section{Evaluation and Results}

\subsection{Evaluation methodology}

The project's standard evaluation is a synthetic fixture harness. It renders a
set of known melodies as audio, runs the pipeline once per fixture, and
compares the output against the known reference. The fixtures cover eight
melodies: a C major scale, an A minor scale, an arpeggio, repeated notes, a
melody with silence, octave leaps, a Twinkle excerpt, and a mixed-rhythm
phrase.

The harness scores two different things, because they measure different
failures:

\begin{itemize}
    \item \textbf{Note detection (pitch and onset F1)}, computed with
    mir\_eval~\cite{mireval}. A predicted note matches a reference note when its
    pitch is correct and its onset is within a tolerance; durations are ignored
    and no offset tolerance is applied. This is run in two flavours: a
    \emph{clean} take that the regression gate expects to pass at F1 $=1.0$, and
    a \emph{realistic} take with wide vibrato, tremolo, pitch drift, jittered
    consonant-gap lengths, and partial consonant closures.
    \item \textbf{Rhythm}, added later, because note F1 is rhythm-blind. This
    scores the quantised onset and duration of each note against an intended
    grid and is described in Section~\ref{sec:results-rhythm}.
\end{itemize}

\subsection{The ground-truth problem}
\label{sec:groundtruth}

Before the results, their central limitation must be stated plainly. \textbf{All
of the ground truth used in this report is synthetic.} The fixtures are
melodies the project itself renders to audio, so their intended pitch, onset,
and duration are known exactly by construction. This is a genuine strength for
regression testing, but it is a genuine weakness for validity: a synthetic hum
is more regular than a real one, and a benchmark built only on synthetic hums
can be saturated while the system still fails on real input.

The project has no labelled corpus of real hummed melodies. The reason is
specific to this task and is the same kind of problem that limits any
transcription benchmark: obtaining ground truth for a real hum requires knowing
the exact pitch, onset, and duration of every note the person intended, and a
person cannot reliably label the pitch of their own hum by ear. As a result,
every note-detection and rhythm figure in this report measures agreement with
\emph{synthetic} ground truth and must be read as a regression signal, not as a
measure of real-world accuracy. Section~\ref{sec:eval-limits} proposes a
concrete remedy.

\subsection{Note-detection results}
\label{sec:results-note}

On the clean fixtures the pipeline scores a mean note F1 of $1.000$. The more
informative number is the realistic take, where the segmenter historically
struggled. Table~\ref{tab:note-progression} records the mean realistic F1
across three engineering iterations. The first raised it from $0.799$ to
$0.931$ by making segmentation and consolidation grid-aware; the second raised
it to $1.000$ by adding the spectral octave correction of
Section~\ref{sec:octave}, which fixed the last failing fixture (the octave
leaps, previously returned an octave flat).

\begin{table}[H]
\centering
\caption{Mean note F1 on the realistic fixtures across three iterations. The
realistic take is the expressive one (wide vibrato, tremolo, drift, jittered
consonant gaps, partial closures).}
\label{tab:note-progression}
\begin{tabular}{l l l}
\toprule
\textbf{Iteration} & \textbf{Change} & \textbf{Mean realistic F1} \\
\midrule
Baseline              & coarse valley splitter                 & $0.799$ \\
Grid-aware            & grid-aware segment + consolidate        & $0.931$ \\
Octave correction     & spectral subharmonic repair             & $1.000$ \\
\bottomrule
\end{tabular}
\end{table}

Table~\ref{tab:note-perfixture} gives the per-fixture note metrics on the
realistic take at the current state. Every fixture scores $1.000$ on precision,
recall, and F1, and the key is detected correctly in every case. This is the
saturation referred to throughout this report: the note-detection half of the
synthetic harness no longer discriminates between a better and a worse
segmenter, so it can no longer drive improvement on its own.

\begin{table}[H]
\centering
\caption{Per-fixture note metrics on the realistic take, current state. P,
R, F1 are precision, recall, and F1 from mir\_eval; ref/est is
reference versus detected note counts.}
\label{tab:note-perfixture}
\begin{tabular}{l l l l l l}
\toprule
\textbf{Fixture} & \textbf{P} & \textbf{R} & \textbf{F1} & \textbf{Key} &
\textbf{ref/est} \\
\midrule
c\_major\_scale  & $1.00$ & $1.00$ & $1.00$ & C major & 8/8 \\
a\_minor\_scale  & $1.00$ & $1.00$ & $1.00$ & A minor & 8/8 \\
arpeggio         & $1.00$ & $1.00$ & $1.00$ & C major & 7/7 \\
repeated\_notes  & $1.00$ & $1.00$ & $1.00$ & C major & 5/5 \\
with\_silence    & $1.00$ & $1.00$ & $1.00$ & C major & 5/5 \\
octave\_leaps    & $1.00$ & $1.00$ & $1.00$ & C major & 5/5 \\
twinkle          & $1.00$ & $1.00$ & $1.00$ & C major & 14/14 \\
mixed\_rhythm    & $1.00$ & $1.00$ & $1.00$ & C major & 5/5 \\
\midrule
\textbf{mean}    &        &        & $1.000$ &        & \\
\bottomrule
\end{tabular}
\end{table}

\subsection{The rhythm evaluation}
\label{sec:results-rhythm}

The note F1 above is rhythm-blind by construction: it ignores durations and
accepts any onset within a tolerance, so it says nothing about whether the
written rhythm is correct. This matters because the two failures reported most
often on real hums are wrong rhythm (durations and onsets on the wrong beats)
and over-splitting (one note returned as tied slivers), and neither is visible
to note F1. A separate rhythm metric was therefore added. It compares the
quantised onset and duration of each note against an \emph{intended grid} read
straight from each fixture's specification, and reports the fraction of notes
with the right onset, the right duration, and both. The headline number is
``both'': right beat and right written length.

Two conditions were run, chosen to match the two things that actually break
real hums. Table~\ref{tab:rhythm} gives the per-fixture results.

\begin{table}[H]
\centering
\caption{Rhythm accuracy under two conditions. ``on'', ``dur'', and ``both''
are the fraction of notes with the correct onset, correct duration, and both;
``err'' is the mean onset error in milliseconds. The jitter condition applies
human timing jitter ($\pm 30$ ms) at the correct tempo; the wrong-tempo
condition applies a $+5\%$ tempo error to on-grid timing.}
\label{tab:rhythm}
\begin{tabular}{l r r r r @{\hskip 1.5em} r r r r}
\toprule
 & \multicolumn{4}{c}{\textbf{Timing jitter ($\pm 30$ ms)}}
 & \multicolumn{4}{c}{\textbf{Wrong tempo ($+5\%$)}} \\
\cmidrule(lr){2-5}\cmidrule(lr){6-9}
\textbf{Fixture} & \textbf{on} & \textbf{dur} & \textbf{both} & \textbf{err} &
\textbf{on} & \textbf{dur} & \textbf{both} & \textbf{err} \\
\midrule
c\_major\_scale  & $1.00$ & $1.00$ & $1.00$ & $0$   & $0.50$ & $1.00$ & $0.50$ & $75$  \\
a\_minor\_scale  & $1.00$ & $1.00$ & $1.00$ & $0$   & $0.50$ & $1.00$ & $0.50$ & $75$  \\
arpeggio         & $1.00$ & $1.00$ & $1.00$ & $0$   & $0.57$ & $1.00$ & $0.57$ & $58$  \\
repeated\_notes  & $1.00$ & $1.00$ & $1.00$ & $0$   & $1.00$ & $0.67$ & $0.67$ & $0$   \\
with\_silence    & $1.00$ & $0.80$ & $0.80$ & $0$   & $0.40$ & $1.00$ & $0.40$ & $90$  \\
octave\_leaps    & $1.00$ & $1.00$ & $1.00$ & $0$   & $1.00$ & $1.00$ & $1.00$ & $0$   \\
twinkle          & $0.77$ & $0.85$ & $0.69$ & $126$ & $0.27$ & $0.64$ & $0.18$ & $855$ \\
mixed\_rhythm    & $1.00$ & $1.00$ & $1.00$ & $0$   & $1.00$ & $1.00$ & $1.00$ & $0$   \\
\midrule
\textbf{mean}    &        &        & $\mathbf{0.937}$ & & & & $\mathbf{0.602}$ & \\
\bottomrule
\end{tabular}
\end{table}

Two findings stand out. First, on-grid at the correct tempo the quantiser is
essentially perfect (a separate on-grid run, not shown, scores $1.000$), so the
quantisation logic itself is sound. Under human timing jitter the mean falls
only to $0.937$: most fixtures stay perfect, but the longest fixture, Twinkle,
drops to $0.69$ and loses a note, so drift bites longer pieces even at the right
tempo. Second, and more importantly, a $5\%$ tempo error collapses the mean to
$0.602$. Scales fall to $0.50$, the silence fixture to $0.40$, Twinkle craters
to $0.18$ with a mean onset error of $855$ ms, and the repeated-notes fixture
even loses notes ($5$ detected as $3$), because a wrong tempo misfires the
grid-aware segmenter as well as the quantiser.

Figure~\ref{fig:rhythm} plots the two conditions side by side.

\begin{figure}[H]
\centering
\reportfig{0.85}{eval-rhythm.png}{Rhythm ``both'' accuracy per fixture under
human timing jitter (correct tempo) and under a $5\%$ tempo error. The tempo
error, not the jitter, is what collapses the score.}
\caption{Rhythm accuracy under timing jitter versus a wrong tempo.}
\label{fig:rhythm}
\end{figure}

\subsection{Interpreting the scores}
\label{sec:interpret}

The note-detection result and the rhythm result must be read together and read
carefully. A saturated note F1 of $1.000$ does not mean the system transcribes
real hums correctly; it means the synthetic note-detection benchmark can no
longer tell a good segmenter from a bad one. The honest reading is that the
project has exhausted what its current synthetic note benchmark can teach, not
that the transcription problem is solved. Informal use on real hums continues to
fail, and there is no measurement in this report that contradicts that, because
there is no real ground truth to measure against.

The rhythm evaluation is more useful precisely because it is not yet saturated.
It localises the dominant real-world failure to a single cause: sensitivity to
tempo error. A $5\%$ tempo mismatch, which is well within the error of the
system's own tempo detector on real input, is enough to drive onsets off the
integer beats, and onsets off the integer beats are what notation renders as
tied slivers. This ties the user-visible ``wrong rhythm'' and ``over-split''
complaints to the same root cause and points directly at the tempo component as
the highest-value target.

\subsection{The tempo-detection weakness}
\label{sec:bpm}

The rhythm result above is a controlled injection of a tempo error. In real use
the error is supplied by the system's own tempo detector, and it is worth
stating why that detector is fragile, because it is the binding constraint on
real-world rhythm accuracy.

As described in Section~\ref{sec:tempo-design}, the detector returns a single
global tempo: it takes the autocorrelation of the onset-strength envelope over
the whole clip, weights it by a prior centred near a comfortable humming tempo,
and returns the dominant period. This has three consequences on real input.
First, it assumes a constant tempo; a hum that speeds up or slows down has onset
energy at more than one period, and the autocorrelation returns a blurry
compromise rather than either true tempo. Second, the prior actively biases the
estimate toward the centre of its range, so a genuinely slow hum is pulled
upward. Third, it relies on clean, evenly spaced onsets, which real legato or
noisy humming does not always provide. In informal testing, a hum of a roughly
$76$ beat-per-minute melody that included a faster passage was reported as about
$101$, a value suspiciously close to the prior, which is the signature of the
prior dominating a weak autocorrelation. This weakness is the leading item of
future work in Section~\ref{sec:directions}.

% =========================================================================
% SECTION 5: LIMITATIONS AND FUTURE WORK
% =========================================================================
\section{Limitations and Future Work}

\subsection{Evaluation limitations}
\label{sec:eval-limits}

The main limitation is the one stated in Section~\ref{sec:groundtruth}: the
evaluation is entirely synthetic, and no labelled corpus of real hummed
melodies exists. Two concrete consequences follow. The note-detection benchmark
is saturated and can no longer drive improvement. And the over-splitting
failure, one of the two dominant real-world complaints, is still not measured
at all, because synthetic vibrato and tremolo are too regular to produce the
spurious narrow energy dips that real shimmer, breath, and creak produce.

The recommended remedy is not a larger synthetic set but a small set of real,
labelled hums. The obstacle is ground truth, since a person cannot label the
pitch of their own hum by ear. The plan, refined during this project, inverts
the labelling. Rather than hum first and label afterward, the label is fixed
before the recording: the user hums a melody they know, to a click at a known
tempo, and the rhythm ground truth comes from the known tune and the click
grid. The pitch of each note is confirmed with an external pitch-finder
application used as an assistive aid, with the user verifying each note against
the tune they knowingly hummed, rather than trusting the external tool as an
oracle. This is important because an external pitch finder is itself another
pitch detector; if it is trusted blindly and it happens to share an error with
the pipeline, the evaluation would score that shared error as correct. Keeping
the user in the loop, and drawing the rhythm truth from the click rather than
from the tool, avoids that circularity. The resulting verified recordings, with
their known notes and tempo, would give the project its first real ground truth.

\subsection{System deficiencies}

Several deficiencies of the current system are known and tracked.
Table~\ref{tab:deficiencies} lists them together with the planned remedy for
each.

\begin{longtable}{>{\raggedright\arraybackslash}p{0.30\textwidth} >{\raggedright\arraybackslash}p{0.58\textwidth}}
\caption{Known deficiencies of the current system and planned work.}
\label{tab:deficiencies} \\
\toprule
\textbf{Area} & \textbf{Description and planned remedy} \\
\midrule
\endfirsthead
\toprule
\textbf{Area} & \textbf{Description and planned remedy} \\
\midrule
\endhead
\midrule
\multicolumn{2}{r}{\footnotesize\itshape Continued on the next page.} \\
\endfoot
\bottomrule
\endlastfoot
Tempo detection fragility & The ``find my tempo'' estimator returns a single
global tempo biased by a prior, and is unreliable on real, variable-tempo,
slow, or legato humming (Section~\ref{sec:bpm}). Planned remedy: flatten or
remove the prior, and build the onset envelope from the pipeline's own voicing
onsets rather than raw spectral flux. \\
\addlinespace
Over-splitting unmeasured & The over-split failure (one note returned as tied
slivers or repeats) is not exposed by the synthetic harness because synthetic
vibrato and tremolo are too regular. Planned remedy: add amplitude shimmer,
breath, and creak to the synthesiser so spurious narrow dips appear. \\
\addlinespace
No real ground truth & All evaluation is synthetic
(Section~\ref{sec:groundtruth}); real-world accuracy is unmeasured. Planned
remedy: the hum-to-a-known-score capture flow of
Section~\ref{sec:eval-limits}. \\
\addlinespace
Saturated note benchmark & The note-detection F1 is at $1.000$ on all fixtures
and no longer discriminates. Planned remedy: harder fixtures and, above all,
real recordings. \\
\addlinespace
Monophonic only & The transcription pipeline assumes a single voice, one note
at a time. Polyphonic or chordal input is out of scope; the Pitch Finder tab
uses a separate chroma-based method for full songs. \\
\addlinespace
Optional backend availability & The most precise backends (PESTO, FCNF0++) are
optional installs with heavier dependency trees, and FCNF0++ downloads weights
on first use. The pipeline falls back to pyin when they are absent, at a
precision cost. \\
\addlinespace
Fixed metre assumption & Quantisation assumes a simple metre and a known,
steady tempo. Time-signature changes and expressive rubato are not modelled. \\
\end{longtable}

\subsection{Future directions}
\label{sec:directions}

Three directions are recorded for the work after this report, in priority
order.

\begin{itemize}
    \item \textbf{(A) Tempo robustness (leading candidate).} The tempo detector
    is the binding constraint on real-world rhythm accuracy
    (Section~\ref{sec:bpm}), and it is directly testable because the fixtures
    have known tempos and the wrong-tempo rhythm pass is already the dashboard
    for it. Two cheap, high-impact changes: flatten or remove the tempo prior
    so slow tempos are not pulled upward, and build the onset envelope from the
    pipeline's own consonant onsets rather than from raw spectral flux.
    \item \textbf{(B) Real-hum ground truth.} The ultimate unlock, blocked only
    on the labelling method, which Section~\ref{sec:eval-limits} resolves. This
    would let the project measure real-world accuracy for the first time and
    de-saturate the benchmark.
    \item \textbf{(C) Over-split synthesis.} Add amplitude shimmer, breath, and
    creak to the synthesiser so that the over-splitting failure becomes
    visible in the harness and can be regressed against.
\end{itemize}

\subsection{Towards production readiness}

HumJob is a research and academic project, not a production service. It runs
locally, which is a deliberate design choice, but a wider deployment would
additionally require packaging the optional neural backends so they install
reliably, a broader automated test suite over real recordings, and error
handling for the many ways a real microphone recording can be unusable (clipping,
background noise, too quiet). None of these is conceptually hard, but together
they are a meaningful body of work beyond the scope of this course project.

% =========================================================================
% SECTION 6: CONCLUSION
% =========================================================================
\section{Conclusion}

This project set out to transcribe a hummed melody into notes, key, and chords,
by accepting two small constraints on how the user hums and building the whole
pipeline around them. The result, HumJob, is a local-first pipeline of pure
functions over three data types, with five interchangeable pitch backends, a
grid-aware segmenter and consolidator, a spectral octave-correction stage, and a
web application with an interactive editor, a pitch finder, a real-time vocal
trainer, a transposer, and a sing-along trainer with optional, opt-in AI
coaching.

With respect to the research questions, the honest verdicts are mixed.
\textbf{RQ1} is supported by construction and by the evaluation: the two
interface constraints do turn the hard parts into tractable ones, and the
grid-aware stages that exploit them are what carried the realistic note F1 from
$0.799$ to $1.000$. \textbf{RQ2} is supported: the diagnosed failures were
segmentation and quantisation failures over an essentially correct pitch
contour, not pitch-model failures. The clearest case is the octave error, which
was proven to be a Viterbi decoding artifact rather than a fault in the audio,
and was repaired by a spectral test rather than by changing the model.
\textbf{RQ3} produced the most important lesson of the project, and it is a
negative one. The synthetic benchmark the system was tuned against is now
saturated at F1 $=1.000$ while the system is still not reliable on real hums,
because a synthetic hum is more regular than a real one and no labelled real
corpus exists to measure against. The rhythm evaluation, which is not yet
saturated, localises the dominant real-world failure to tempo sensitivity: the
quantiser is perfect on-grid but collapses under a small tempo error, and the
tempo detector that should prevent that error is a single global-tempo
estimator that is fragile on real humming.

The main deliverable is therefore not a headline accuracy number. It is a
working, well-structured pipeline together with an honest evaluation that says
clearly where it stands: strong on a synthetic benchmark that has stopped being
informative, and not yet validated on real input. The most valuable next step
is to build a small, trustworthy set of real, labelled hums, so that the
system's real-world accuracy can be measured directly rather than inferred from
a saturated synthetic proxy.

% =========================================================================
% REFERENCES
% =========================================================================
\newpage
\addcontentsline{toc}{section}{References}
\begin{thebibliography}{99}

\bibitem{pyin} M. Mauch and S. Dixon. \emph{pYIN: A fundamental frequency
estimator using probabilistic threshold distributions.} IEEE International
Conference on Acoustics, Speech and Signal Processing (ICASSP), 2014.

\bibitem{crepe} J. W. Kim, J. Salamon, P. Li, and J. P. Bello. \emph{CREPE: A
Convolutional Representation for Pitch Estimation.} ICASSP, 2018.
arXiv:1802.06182. Implementation: torchcrepe,
\url{https://github.com/maxrmorrison/torchcrepe}.

\bibitem{pesto} A. Riou, S. Lattner, G. Hadjeres, and G. Peeters. \emph{PESTO:
Pitch Estimation with Self-supervised Transposition-equivariant Objective.}
International Society for Music Information Retrieval Conference (ISMIR), 2023.
\url{https://github.com/SonyCSLParis/pesto}.

\bibitem{penn} M. Morrison, C. Hsieh, N. Pruyne, and B. Pardo.
\emph{Cross-domain Neural Pitch and Periodicity Estimation.} 2023.
arXiv:2301.12258. Implementation: penn,
\url{https://github.com/interactiveaudiolab/penn}.

\bibitem{basicpitch} R. M. Bittner, J. J. Bosch, D. Rubinstein,
G. Meseguer-Brocal, and S. Ewert. \emph{A Lightweight Instrument-Agnostic Model
for Polyphonic Note Transcription and Multipitch Estimation.} ICASSP, 2022.
\url{https://github.com/spotify/basic-pitch}.

\bibitem{krumhansl} C. L. Krumhansl. \emph{Cognitive Foundations of Musical
Pitch.} Oxford University Press, 1990. (Krumhansl-Kessler / Krumhansl-Schmuckler
probe-tone key profiles.)

\bibitem{librosa} B. McFee, C. Raffel, D. Liang, D. P. W. Ellis, M. McVicar,
E. Battenberg, and O. Nieto. \emph{librosa: Audio and Music Signal Analysis in
Python.} Proceedings of the 14th Python in Science Conference (SciPy), 2015.
\url{https://librosa.org/}.

\bibitem{music21} M. S. Cuthbert and C. Ariza. \emph{music21: A Toolkit for
Computer-Aided Musicology and Symbolic Music Data.} ISMIR, 2010.
\url{https://web.mit.edu/music21/}.

\bibitem{verovio} L. Pugin, R. Zitellini, and P. Roland. \emph{Verovio: A
library for Engraving MEI Music Notation into SVG.} ISMIR, 2014.
\url{https://www.verovio.org/}.

\bibitem{mireval} C. Raffel, B. McFee, E. J. Humphrey, J. Salamon, O. Nieto,
D. Liang, and D. P. W. Ellis. \emph{mir\_eval: A Transparent Implementation of
Common MIR Metrics.} ISMIR, 2014. \url{https://github.com/craffel/mir_eval}.

\end{thebibliography}

% =========================================================================
% APPENDIX A - SYSTEM DEMONSTRATION (SCREENSHOTS)
% =========================================================================
\newpage
\appendix
\addcontentsline{toc}{section}{Appendix A: System demonstration}
\section*{Appendix A: System demonstration}
\renewcommand{\thesubsection}{A.\arabic{subsection}}
\setcounter{subsection}{0}

This appendix shows the running web application from a user's perspective. Each
subsection corresponds to one tab. The screenshots are captured against a local
instance of the stack; small visual details depend on the user's browser and
operating system.

\subsection{Transcriber: Auto mode}

The Transcriber tab records or uploads a hum and returns the automatic draft:
the engraved sheet music, the detected key, and the suggested chords, with
downloads for MIDI and MusicXML.

\begin{figure}[H]
\centering
\reportfig{0.85}{screenshot-transcriber.png}{The Transcriber tab in Auto mode:
the engraved server-rendered sheet, the detected key, and the export buttons.}
\caption{Transcriber: Auto mode.}
\label{fig:scr-transcriber}
\end{figure}

\subsection{Transcriber: Manual mode}

Manual mode is the in-browser staff editor for correcting the draft. The
reference pitch strip shows the hummed contour against the chosen notes, and the
editing toolbar offers pitch, duration, merge, split, delete, and insert with
undo and redo.

\begin{figure}[H]
\centering
\reportfig{0.85}{screenshot-manual.png}{The Transcriber tab in Manual mode: the
editable staff, the reference pitch strip, and the editing toolbar.}
\caption{Transcriber: Manual mode.}
\label{fig:scr-manual}
\end{figure}

\subsection{Pitch Finder}

The Pitch Finder tab analyses an uploaded clip for key, tempo, and a Camelot
code, using the chroma-based method that works on polyphonic audio.

\begin{figure}[H]
\centering
\reportfig{0.85}{screenshot-pitchfinder.png}{The Pitch Finder tab: detected
key, tempo, Camelot code, and technical statistics for an uploaded clip.}
\caption{Pitch Finder.}
\label{fig:scr-pitchfinder}
\end{figure}

\subsection{Realtime and vocal trainer}

The Realtime tab is a client-side pitch monitor and guitar tuner with
vocal-training drills: a target-note lane, a match game, a scale trainer,
vibrato analysis, a vocal range finder, and a per-session progress panel.

\begin{figure}[H]
\centering
\reportfig{0.85}{screenshot-realtime.png}{The Realtime tab: the live voice
monitor with the practice sub-mode selector (Free / Match game / Scale trainer
/ Range), the tuner meter, the steadiness readout, and the live pitch graph.}
\caption{Realtime voice monitor and vocal trainer.}
\label{fig:scr-realtime}
\end{figure}

\subsection{Transposer}

The Transposer tab shifts a score to a new key. File mode transposes an
uploaded MIDI or MusicXML file server-side and offers Camelot-compatible key
presets; hum mode transposes the last transcription client-side.

\begin{figure}[H]
\centering
\reportfig{0.85}{screenshot-transposer.png}{The Transposer tab in hum mode: the
just-hummed melody transposed client-side, with the shift control, the To-key
picker, and the transposed chords. Chord functions (the Roman numerals) are
preserved, so only the chord names move.}
\caption{Transposer: hum mode.}
\label{fig:scr-transposer}
\end{figure}

\subsection{Sing-Along and AI coaching}

The Sing-Along tab plays an uploaded score as a guide, tracks the voice live,
and scores each note. After a take it shows a deterministic analysis of the
performance offline, and, on request, sends a numeric summary to an external
language model for spoken-language coaching (the one non-local, opt-in feature).

\begin{figure}[H]
\centering
\reportfig{0.85}{screenshot-singalong.png}{The Sing-Along tab after a take: the
per-note scoring overview, the deterministic analysis (pitch bias, drift, leap
versus step accuracy, weakest register, worst notes), and the optional
``Get coaching'' panel with its language selector.}
\caption{Sing-Along and AI coaching.}
\label{fig:scr-singalong}
\end{figure}

\end{document}
